% Generated draft: requires mathematical and source review.
\documentclass[11pt]{article}
\usepackage[margin=25mm]{geometry}
\usepackage{fontspec}
\setmainfont{Latin Modern Roman}
\usepackage{amsmath,amssymb,mathtools,unicode-math}
\setmathfont{Latin Modern Math}
\usepackage{enumitem,needspace,xurl,hyperref}
\setlength{\parindent}{0pt}
\setlength{\parskip}{5pt}
\setlength{\emergencystretch}{3em}
\newcommand{\N}{\mathbb{N}}
\newcommand{\Z}{\mathbb{Z}}
\newcommand{\Q}{\mathbb{Q}}
\newcommand{\R}{\mathbb{R}}
\newcommand{\C}{\mathbb{C}}
\newcommand{\F}{\mathbb{F}}
\newcommand{\A}{\mathbb{A}}
\newcommand{\PP}{\mathbb{P}}
\newcommand{\E}{\mathbb{E}}
\newcommand{\eps}{\varepsilon}
\newcommand{\dd}{\,\mathrm{d}}
\begin{document}
\section*{UnsolvedMath ID: 30006326}
\textbf{Source problem code:} catalogue entry
\subsection{Definitions and mathematical statement}
The supplied record concerns the Whittaker special case of a global geometric duality conjecture. Its intended assertion is incomplete; the following records the stated objects and the requested identification without supplying missing hypotheses or a missing spectral object.

Let \(G\) and \(U\subseteq G\) be the group data of the source, and let \(\mathbb G_a\) denote the additive group. Given a character
\[
\psi\colon U\longrightarrow\mathbb G_a
\]
described in the source as generic, put \(X=G/U\). Let \(\Psi\to X\) be the associated \(\mathbb G_a\)-bundle. The source defines the Whittaker period sheaf by
\[
\mathcal W:=P_{G/U,\Psi}.
\]
The ground field, hypotheses on \(G\) and \(U\), meaning of genericity, construction of \(P_{G/U,\Psi}\), and category containing \(\mathcal W\) are not specified in the supplied fragment.

Write \(G^\vee\) for the dual group appearing in the source and \(\operatorname{Loc}_{G^\vee}\) for its moduli object of local systems. For a \(G^\vee\)-local system \(\rho\in\operatorname{Loc}_{G^\vee}\), the source supposes that \(\mathcal F_\rho\) is an associated Hecke eigensheaf and states that its image under the global geometric Langlands correspondence is
\[
\delta_\rho:=(i_\rho)_*\mathcal O_{\mathrm{pt}}
\quad\text{in}\quad
\operatorname{QC}^{!}(\operatorname{Loc}_{G^\vee}),
\]
where \(i_\rho\colon\mathrm{pt}\to\operatorname{Loc}_{G^\vee}\) denotes the point represented by \(\rho\). This is the correspondence stipulated by the source; the precise moduli convention, sheaf theory, and functor are not supplied.

\textbf{Requested identification.} In the intended global geometric Langlands setting, identify the image of \(\mathcal W\) with the corresponding Whittaker spectral object. The supplied fragment does not define that spectral object or state the encompassing duality conjecture, so this request cannot be expressed here as an isomorphism between two independently specified objects.

The fragment further describes a favorable situation in which \(\rho\) is a smooth point of \(\operatorname{Loc}_{G^\vee}\) and a dual space \(X^\vee\) has an isolated fixed point \(x\). It says that an unspecified Hom space can then be understood as a suitably twisted and sheared symmetric algebra on
\[
H^1\!\left(\Sigma;(T_xX^\vee)_\rho\right).
\]
Here \(\Sigma\), \(X^\vee\), the action defining the fixed point, the representation needed to form \((T_xX^\vee)_\rho\), the two objects in the Hom space, and the twisting and shearing conventions are absent. Consequently, this description is retained as a source qualification, not asserted as a precise Hom-space identity. The extracted text writes \(\{x\}=X^\vee\) while referring to an isolated fixed point; whether a fixed-point locus was intended is unresolved. The source interprets the period sheaf as a geometric avatar of a period integral and the cohomological description as a geometric avatar of a corresponding \(L\)-value. It also says that its favorable conditions correspond to cases where the \(L\)-function has no poles, but neither the \(L\)-function nor the required evaluation point is specified.
\subsection{Short English statement}
The intended problem is to identify the spectral counterpart of the Whittaker period sheaf under global geometric Langlands. The supplied excerpt defines the period sheaf and sketches a favorable cohomological description, but omits the spectral object and essential hypotheses. It therefore supports an incomplete conjectural target requiring mathematical review.
\subsection{Sources}
\begin{itemize}
\item Local UnsolvedMath registry, \texttt{data/unsolvedmath/problems.json}, supplied record with permanent ID \(30006326\), title \emph{Whittaker Special Cases of Global Geometric Duality}, exact source problem code \texttt{catalogue entry}. The record gives proposed year 2025 and names no individual proposer.
\item \emph{Arbeitsgemeinschaft: Relative Langlands Duality} (2025), Oberwolfach Reports, pp.\ 811--882, DOI 10.4171/owr/2025/18. The extracted passage identifies report p.\ 871 and subsection 2.1, ``The Whittaker Case.'' Supplied source URL and literature-source link: \url{https://doi.org/10.4171/owr/2025/18}.
\item \emph{Relative Langlands Duality}, supplied literature-source link: \url{https://arxiv.org/abs/2409.04677}. No authors or further bibliographic details are supplied in the record.
\item Literature assessment dated exactly \texttt{2026-08-22}: the supplied curation labels the item open, explains that the fragment does not state a standalone theorem, and attributes that classification to the unresolved encompassing global duality claim. This is supplied dated curation, not an independently verified assessment. The background expressly requests independent verification of time-sensitive status and cited claims.
\item The registry supplies the label \texttt{corrected\_verified} and a verification note claiming that the official report was checked and its notation repaired. These are source-record assertions; no independent source consultation or verification is claimed in this draft.
\end{itemize}
\end{document}
